\section{Tres tipos de Transformer: a partir de la base común del sequence model}

El capítulo anterior discutió por qué conviene estudiar la estructura; este capítulo primero la explica con claridad. Encoder-decoder, encoder-only y decoder-only no son tres invenciones sin relación entre sí: las tres convierten una entrada discreta en una secuencia de vectores y después usan attention para modelar la interacción entre los elementos. Las diferencias provienen principalmente de si los parámetros se comparten, de cómo la attention mask limita el alcance visible y de cómo se entrena y se usa la salida.

\subsection{La interfaz de tarea de las tres arquitecturas}

Esta sección presenta primero el panorama general y después entra paso a paso en las figuras. Encoder-decoder asigna dos stacks distintos a la entrada y a la salida; encoder-only comprime la entrada en una representación y por lo general requiere un task-specific head; decoder-only concatena la entrada y el objetivo, y genera token por token en un único causal stack. Cada interfaz implica supuestos distintos sobre la tarea.

\lecturefigure{hyung-slide-017.jpg}{Los tres tipos de arquitectura Transformer: encoder-decoder, encoder-only y decoder-only}{Hyung Won Chung official deck, p.~17}

\begin{knowledgebox}{Glosario: qué problema resuelve cada arquitectura}
\begin{tabularx}{\textwidth}{L{0.22\textwidth} X}
\toprule
Arquitectura & Mecanismo central y relación con el curso \\
\midrule
Encoder-decoder & El encoder codifica la entrada de forma bidireccional; el decoder genera el objetivo mediante causal self-attention y cross-attention. Es adecuada para sequence-to-sequence explícito. \\
Encoder-only & Interacción bidireccional sobre toda la secuencia; la salida suele usarse para clasificación, recuperación o token labeling. Para generar una secuencia necesita un mecanismo de generación adicional. \\
Decoder-only & La entrada y la salida previa están en la misma secuencia causal; comparte parámetros y tiene una interfaz unificada, por lo que es adecuada de forma natural para autoregressive generation y para la continuación en varios turnos. \\
\bottomrule
\end{tabularx}
\end{knowledgebox}

\subsection{De los caracteres a la secuencia de vectores, y de ahí a la interaction}

Antes de clasificar las arquitecturas hay que entender el pipeline común. Primero, el tokenizer asigna al texto token IDs; después, la embedding table asigna a cada ID un vector de dimensión $d_{\text{model}}$; por último, el sequence model permite que las distintas posiciones intercambien información según el attention pattern. Lo que el modelo procesa es una secuencia de vectores, no directamente el «significado de las palabras».

\lecturefigure{hyung-slide-018.jpg}{La secuencia original de caracteres es el punto de partida del pipeline del Transformer}{Hyung Won Chung official deck, p.~18}

\lecturefigure{hyung-slide-021.jpg}{Pipeline completo: tokenización, embedding y sequence model}{Hyung Won Chung official deck, p.~21}

Si el token ID es $i_t$ y la embedding matrix es $E\in\mathbb{R}^{V\times d}$, entonces
\[
h_t^{(0)}=E[i_t]+p_t.
\]
Aquí $V$ es el tamaño del vocabulario, $d$ es la hidden dimension, $E[i_t]$ es el vector del token número $i_t$, $p_t$ es la representación de posición y $h_t^{(0)}$ es el estado que entra a la primera capa. La tokenización determina la longitud de la secuencia y los límites de la segmentación; el embedding determina la representación continua inicial. Ambos influyen en el costo de los cálculos posteriores.

\teachervoice{Entre 00:16:11 y 00:17:26, Hyung Won define primero el Transformer como un sequence model: los elementos de entrada pueden ser palabras, imágenes u otros objetos, y lo central es modelar la interaction entre ellos. Este orden evita separar la fórmula de attention de la interfaz de datos.}

\subsection{Encoder-decoder: tres attention roles}

Ahora pasamos al Transformer original. Al leer la figura, conviene separar primero los tres flujos de información: la self-attention bidireccional dentro del encoder, la self-attention causal dentro del decoder y la cross-attention con la que el decoder consulta la representación del encoder. Cada arista corresponde a un permiso del tipo «qué posición puede leer qué posición».

\lecturefigure{hyung-slide-022.jpg}{Arquitectura encoder-decoder completa y sus tres attention roles}{Hyung Won Chung official deck, p.~22}

\begin{importantbox}{Self-attention y cross-attention}
Dados query $Q$, key $K$ y value $V$, la scaled dot-product attention es
\[
\operatorname{Attn}(Q,K,V)=\operatorname{softmax}\!\left(\frac{QK^{\top}}{\sqrt{d_k}}+M\right)V.
\]
Aquí $d_k$ es la dimensión de key y $M$ es la mask de visibilidad. En la self-attention, $Q,K,V$ provienen de la misma secuencia; en la cross-attention, $Q$ proviene del decoder y $K,V$ provienen de la salida del encoder.
\end{importantbox}

\lecturefigure{hyung-slide-023.jpg}{El encoder usa bidirectional self-attention para modelar la entrada completa}{Hyung Won Chung official deck, p.~23}

La bidirectional attention (atención bidireccional) permite que cualquier posición de la entrada lea cualquier otra posición. Si la longitud de la secuencia es $n$, la mask de visibilidad puede escribirse como $M_{ij}=0$; no hay enmascaramiento del futuro. Es adecuada para una entrada que se proporciona de una sola vez, porque todos los tokens ya se conocen al codificar, pero también implica que un cambio al final de la entrada afecta la representación de todas las posiciones anteriores.

\lecturefigure{hyung-slide-025.jpg}{El decoder usa causal self-attention y solo lee la posición actual y el historial a su izquierda}{Hyung Won Chung official deck, p.~25}

La causal attention (atención causal) usa una mask triangular superior para impedir la lectura del futuro:
\[
M_{ij}=\begin{cases}
0,&j\le i,\\
-\infty,&j>i.
\end{cases}
\]
Aquí $i$ es la posición actual de query y $j$ es la posición de key que se lee. Después de la softmax, el peso de las posiciones futuras se vuelve cero; por eso, durante el entrenamiento pueden calcularse todas las posiciones en paralelo y, durante la inferencia, se mantiene la condición autoregressive.

\lecturefigure{hyung-slide-027.jpg}{La cross-attention conecta las posiciones del objetivo en el decoder con la representación de entrada del encoder}{Hyung Won Chung official deck, p.~27}

\lecturefigure{hyung-slide-028.jpg}{Construcción final del encoder-decoder: codificación bidireccional de la entrada y generación causal del objetivo}{Hyung Won Chung official deck, p.~28}

\begin{knowledgebox}{Lectura de la figura: cómo fluye una ruta de traducción automática}
La entrada ``That is good'' primero interactúa de forma bidireccional dentro del encoder y forma una representación contextualizada de cada token de origen; al generar ``Das ist gut'', el decoder solo puede leer el prefijo del objetivo ya generado y, al mismo tiempo, accede a la última capa del encoder mediante cross-attention. Lo que la figura respalda es la estructura de permisos de información; no indica que los pesos de attention equivalgan por naturaleza a una explicación de alineamiento.
\end{knowledgebox}

\subsection{Diferencias clave entre encoder-only y decoder-only}

La sección anterior usó dos stacks; esta sección examina los dos extremos. Encoder-only produce una representación bidireccional global y, para clasificar, suele leer un token especial y agregar una task-specific layer; decoder-only, en cambio, unifica la entrada y la salida en una sola secuencia y usa el mismo conjunto de parámetros y causal self-attention tanto para comprender el contexto como para generar.

\lecturefigure{hyung-slide-032.jpg}{Forma completa de encoder-only: representación bidireccional, token especial y cabeza de tarea}{Hyung Won Chung official deck, p.~32}

\begin{warningbox}{Encoder-only no significa «no puede producir ninguna secuencia»}
Esta diapositiva se refiere a la interfaz predeterminada: un solo encoder stack no incluye por sí mismo una autoregressive generation path. Es posible agregarle un decoder, un span head, el llenado iterativo de huecos u otro mecanismo de generación; el costo es que el sistema deja de ser la estructura encoder-only mínima de la figura. Lo que el curso compara son los supuestos integrados, no la capacidad absoluta en sentido lógico.
\end{warningbox}

\lecturefigure{hyung-slide-037.jpg}{Forma completa de decoder-only: entrada y objetivo concatenados, parámetros compartidos y causal attention}{Hyung Won Chung official deck, p.~37}

La interfaz unificada de decoder-only puede escribirse como
\[
p(y\mid x)=\prod_{t=1}^{|y|}p(y_t\mid x,y_{<t}).
\]
Aquí $x$ es la secuencia de tokens de entrada, $y$ es la secuencia objetivo y $y_{<t}$ es el prefijo del objetivo ya generado. La entrada y el objetivo usan el mismo stack, pero la loss normalmente puede calcularse solo en las posiciones del objetivo; «compartir parámetros» no significa que el modelo no pueda distinguir los roles, porque la posición, los tokens separadores y el propio contexto siguen proporcionando condiciones.

\subsection{Resumen de la sección}

Los tres tipos de Transformer comparten la base de tokenización, embedding, attention y MLP; las diferencias principales están en la visibilidad de la información, la agrupación de parámetros y la interfaz de tarea. Una vez entendidas estas diferencias, el siguiente capítulo ya no trata encoder-decoder y decoder-only como categorías inconmensurables, sino que transforma de manera continua la primera en la segunda mediante cuatro operaciones, lo que deja al descubierto el verdadero sesgo inductivo.
